% Converted from S6_Concurrency_conflicts/reviewer_response_package_20261116_claude_v3/RESPONSE_R2_17_EN.md (2026-11-17)
We thank the reviewer for this question. In the report-generation workflow of ThyroidXAgent, lymph-node screening and nodule analysis do not run concurrently, and the workflow does not match objects between the two branches. Duplication is avoided by anatomical region: lymph-node screening analyses only neck-region images, and nodules assigned to the neck region are not described as thyroid nodules in the report. We describe this mechanism below and quantified it in the 450 audit runs described in our response to Comment 16.\par

\textbf{Execution order.} Execution can be described at two levels, between tools and within a tool. Between tools, the only parallel step is preprocessing, in which three per-image classifiers (colour-Doppler detection, anatomical region classification and nodule-presence detection) run concurrently, each writing a separate field of the image record. Within a tool, the per-image tools, such as nodule segmentation, feature classification, pixel-spacing prediction and lymph-node screening, take a batch of images in one call and return one result per image. After preprocessing, the tools for different tasks, including nodule segmentation and lymph-node screening, are called one at a time by the executor. In the 450 audit runs, none of the 2,347 tool calls overlapped in time with another.\par

\textbf{Separation by region.} Preprocessing assigns each image to one region: a longitudinal or transverse view of the left or right lobe, a transverse view of both lobes or the isthmus, or the neck. By the definitions of the tools in the toolset, lymph-node screening accepts only neck images, and nodule segmentation analyses the images flagged by the nodule-presence detector, which can include neck images. The fusion tool groups nodules by lobe, so nodules in neck images are never merged with nodules in thyroid images. When the report is assembled, nodules in the neck region are not described as thyroid nodules.\par

\textbf{What happened in the audit runs.} The nodule-presence detector flagged 72 neck images in 27 of the 450 runs (6.0\%; 95\% CI 2.7--10.0\%; 9 cases), and segmentation detected a structure in 69 of them. In 60 of these images (24 runs, 5.3\%; 95\% CI 2.0--9.3\%; 8 cases), lymph-node screening was also positive, that is, both branches detected a structure in the same image. All 69 nodules segmented in neck images were fused as neck-region nodules, none was combined with a nodule from a thyroid image, and none appeared in a nodule description. They also did not change the TI-RADS category stated in the impression (24 runs; KMVE reports do not state a category).\par

The same separation held in the 1,042 test reports of the original study. There, 57 images in 35 reports were processed by both branches, and none of the 57 nodules fused from these images was described as a thyroid nodule.\par
